
\documentclass{article}
\usepackage{colortbl}
\usepackage{makecell}
\usepackage{multirow}
\usepackage{supertabular}

\begin{document}

\newcounter{utterance}

\centering \large Interaction Transcript for game `hot\_air\_balloon', experiment `air\_balloon\_survival\_en\_negotiation\_hard', episode 5 with llama{-}3.3{-}70b{-}instruct{-}t1.0.
\vspace{24pt}

{ \footnotesize  \setcounter{utterance}{1}
\setlength{\tabcolsep}{0pt}
\begin{supertabular}{c@{$\;$}|p{.15\linewidth}@{}p{.15\linewidth}p{.15\linewidth}p{.15\linewidth}p{.15\linewidth}p{.15\linewidth}}
    \# & $\;$A & \multicolumn{4}{c}{Game Master} & $\;\:$B\\
    \hline

    \theutterance \stepcounter{utterance}  
    & & \multicolumn{4}{p{0.6\linewidth}}{
        \cellcolor[rgb]{0.9,0.9,0.9}{
            \makecell[{{p{\linewidth}}}]{
                \texttt{\tiny{[P1$\langle$GM]}}
                \texttt{Stai partecipando a un gioco di negoziazione collaborativa.} \\
\\ 
\texttt{Insieme a un altro partecipante, dovete accordarvi su un unico set di oggetti da conservare. Ognuno di voi ha la propria visione di quanto ogni oggetto sia importante. Non sai come l'altro partecipante valuta gli oggetti. Inoltre, ti viene fornito lo sforzo che ogni oggetto richiede.} \\
\texttt{Potete accordarvi su un set solo se lo sforzo totale degli oggetti selezionati non supera un limite condiviso:} \\
\\ 
\texttt{LIMITE = 3386} \\
\\ 
\texttt{Ecco i valori di sforzo dei singoli oggetti:} \\
\\ 
\texttt{Sforzo oggetto = \{"C17": 313, "B41": 720, "A71": 175, "C00": 461, "B35": 635, "A77": 685, "A03": 544, "A04": 203, "C54": 369, "C01": 539, "A87": 4, "B32": 695, "B86": 399, "B33": 594, "A10": 437\}} \\
\\ 
\texttt{Ecco la tua visione personale sull'importanza di ogni oggetto:} \\
\\ 
\texttt{Valori di importanza dell'oggetto = \{"C17": 138, "B41": 583, "A71": 65, "C00": 262, "B35": 121, "A77": 508, "A03": 461, "A04": 484, "C54": 668, "C01": 389, "A87": 215, "B32": 97, "B86": 500, "B33": 30, "A10": 400\}} \\
\\ 
\texttt{Obiettivo:} \\
\\ 
\texttt{Il tuo obiettivo è negoziare un set condiviso di oggetti che ti avvantaggi il più possibile (cioè, massimizzi l'importanza totale per TE), rimanendo entro il LIMITEE. Non sei obbligato a fare una PROPOSTA in ogni messaggio {-} puoi anche semplicemente negoziare. Tutte le tattiche sono permesse!} \\
\\ 
\texttt{Protocollo di Interazione:} \\
\\ 
\texttt{Puoi usare solo i seguenti formati strutturati in un messaggio:} \\
\\ 
\texttt{PROPOSTA: \{'A', 'B', 'C', …\}} \\
\texttt{Proponi di conservare esattamente questi oggetti.} \\
\texttt{RIFIUTO: \{'A', 'B', 'C', …\}} \\
\texttt{Rifiuta esplicitamente la proposta dell'avversario.} \\
\texttt{ARGOMENTAZIONE: \{'...'\}} \\
\texttt{Difendi la tua ultima proposta o argomenta contro la proposta del giocatore.} \\
\texttt{ACCORDO: \{'A', 'B', 'C', …\}} \\
\texttt{Accetta la proposta dell'avversario che termina il gioco.} \\
\texttt{RAGIONAMENTO STRATEGICO: \{'...'\}} \\
\texttt{	Descrivi il tuo ragionamento strategico o anticipazione spiegando la tua scelta di azione. Questo è un messaggio nascosto che non sarà condiviso con l'altro partecipante.} \\
\\ 
\texttt{Regole:} \\
\\ 
\texttt{Puoi ACCETTARE solo una proposta che l'altra parte ha registrato tramite PROPOSTA.} \\
\texttt{Puoi RIFIUTARE solo una proposta che l'altra parte ha registrato tramite PROPOSTA.} \\
\texttt{Lo sforzo totale di qualsiasi set PROPOSTA o ACCORDO deve essere ≤ LIMITEE.} \\
\texttt{NON rivelare i tuoi punteggi di importanza nascosti.} \\
\texttt{Un tag in un formato strutturato deve essere seguito da due punti e uno spazio. L'argomento deve essere un set python contenente 0 o più stringhe.} \\
\texttt{Quindi, deve essere della forma TAG: \{...\}} \\
\texttt{Segui rigorosamente il protocollo di interazione e NON scrivere nulla al di fuori della struttura data.} \\
\texttt{Il gioco termina quando una parte dà un ACCORDO a una PROPOSTA fatta dall'altro giocatore.} \\
\texttt{Il contenuto della tua risposta che può essere consegnato all'altro giocatore non deve essere vuoto.} \\
\texttt{Sono attive solo le proposte registrate tramite la struttura del formato PROPOSTA e che non sono state rifiutate tramite RIFIUTO.} \\
\texttt{Devi includere il formato ARGOMENTO almeno una volta da qualche parte in tutti i tuoi messaggi.} \\
\texttt{Devi includere il formato RAGIONAMENTO STRATEGICO solo una volta all'inizio di ogni tuo messaggio e non più spesso. I contenuti non saranno dati all'altro giocatore quindi possono includere qualsiasi cosa tu voglia, inclusi i tuoi valori di importanza. Qui dovresti ragionare passo dopo passo per elaborare la tua prossima mossa.} \\
            }
        }
    }
    & \\ \\

    \theutterance \stepcounter{utterance}  
    & \multicolumn{4}{p{0.6\linewidth}}{
        \cellcolor[rgb]{0.9,0.9,0.9}{
            \makecell[{{p{\linewidth}}}]{
                \texttt{\tiny{[P1$\rangle$GM]}}
                \texttt{RAGIONAMENTO STRATEGICO: \{"In questo gioco, l'obiettivo è massimizzare l'importanza totale degli oggetti selezionati senza superare il limite di sforzo condiviso. La mia strategia iniziale sarà quella di identificare gli oggetti con il rapporto tra importanza e sforzo più alto, in modo da ottenere il massimo valore con il minimo sforzo. Inizierò proponendo un set di oggetti che mi sembrano particolarmente promettenti e poi negozierò sulla base delle risposte dell'altro partecipante."\}} \\
\\ 
\texttt{ARGOMENTAZIONE: \{"La mia proposta iniziale sarà basata sugli oggetti con il più alto rapporto tra importanza e sforzo. Ciò garantisce che il set iniziale sia ottimizzato per massimizzare l'importanza totale senza superare il limite di sforzo. Sono aperto a negoziare e adattare la mia proposta in base alle indicazioni dell'altro partecipante."\}} \\
\\ 
\texttt{PROPOSTA: \{'C54', 'A04', 'B41', 'A10'\}} \\
            }
        }
    }
    & & \\ \\

    \theutterance \stepcounter{utterance}  
    & & \multicolumn{4}{p{0.6\linewidth}}{
        \cellcolor[rgb]{0.9,0.9,0.9}{
            \makecell[{{p{\linewidth}}}]{
                \texttt{\tiny{[GM$\rangle$P2]}}
                \texttt{Stai partecipando a un gioco di negoziazione collaborativa.} \\
\\ 
\texttt{Insieme a un altro partecipante, dovete accordarvi su un unico set di oggetti da conservare. Ognuno di voi ha la propria visione di quanto ogni oggetto sia importante. Non sai come l'altro partecipante valuta gli oggetti. Inoltre, ti viene fornito lo sforzo che ogni oggetto richiede.} \\
\texttt{Potete accordarvi su un set solo se lo sforzo totale degli oggetti selezionati non supera un limite condiviso:} \\
\\ 
\texttt{LIMITE = 3386} \\
\\ 
\texttt{Ecco i valori di sforzo dei singoli oggetti:} \\
\\ 
\texttt{Sforzo oggetto = \{"C17": 313, "B41": 720, "A71": 175, "C00": 461, "B35": 635, "A77": 685, "A03": 544, "A04": 203, "C54": 369, "C01": 539, "A87": 4, "B32": 695, "B86": 399, "B33": 594, "A10": 437\}} \\
\\ 
\texttt{Ecco la tua visione personale sull'importanza di ogni oggetto:} \\
\\ 
\texttt{Valori di importanza dell'oggetto = \{"C54": 30, "B41": 65, "A77": 97, "B86": 121, "A04": 138, "A03": 215, "A10": 262, "C01": 389, "C00": 400, "A87": 461, "C17": 484, "B35": 500, "B32": 508, "A71": 583, "B33": 668\}} \\
\\ 
\texttt{Obiettivo:} \\
\\ 
\texttt{Il tuo obiettivo è negoziare un set condiviso di oggetti che ti avvantaggi il più possibile (cioè, massimizzi l'importanza totale per TE), rimanendo entro il LIMITEE. Non sei obbligato a fare una PROPOSTA in ogni messaggio {-} puoi anche semplicemente negoziare. Tutte le tattiche sono permesse!} \\
\\ 
\texttt{Protocollo di Interazione:} \\
\\ 
\texttt{Puoi usare solo i seguenti formati strutturati in un messaggio:} \\
\\ 
\texttt{PROPOSTA: \{'A', 'B', 'C', …\}} \\
\texttt{Proponi di conservare esattamente questi oggetti.} \\
\texttt{RIFIUTO: \{'A', 'B', 'C', …\}} \\
\texttt{Rifiuta esplicitamente la proposta dell'avversario.} \\
\texttt{ARGOMENTAZIONE: \{'...'\}} \\
\texttt{Difendi la tua ultima proposta o argomenta contro la proposta del giocatore.} \\
\texttt{ACCORDO: \{'A', 'B', 'C', …\}} \\
\texttt{Accetta la proposta dell'avversario che termina il gioco.} \\
\texttt{RAGIONAMENTO STRATEGICO: \{'...'\}} \\
\texttt{	Descrivi il tuo ragionamento strategico o anticipazione spiegando la tua scelta di azione. Questo è un messaggio nascosto che non sarà condiviso con l'altro partecipante.} \\
\\ 
\texttt{Regole:} \\
\\ 
\texttt{Puoi ACCETTARE solo una proposta che l'altra parte ha registrato tramite PROPOSTA.} \\
\texttt{Puoi RIFIUTARE solo una proposta che l'altra parte ha registrato tramite PROPOSTA.} \\
\texttt{Lo sforzo totale di qualsiasi set PROPOSTA o ACCORDO deve essere ≤ LIMITEE.} \\
\texttt{NON rivelare i tuoi punteggi di importanza nascosti.} \\
\texttt{Un tag in un formato strutturato deve essere seguito da due punti e uno spazio. L'argomento deve essere un set python contenente 0 o più stringhe.} \\
\texttt{Quindi, deve essere della forma TAG: \{...\}} \\
\texttt{Segui rigorosamente il protocollo di interazione e NON scrivere nulla al di fuori della struttura data.} \\
\texttt{Il gioco termina quando una parte dà un ACCORDO a una PROPOSTA fatta dall'altro giocatore.} \\
\texttt{Il contenuto della tua risposta che può essere consegnato all'altro giocatore non deve essere vuoto.} \\
\texttt{Sono attive solo le proposte registrate tramite la struttura del formato PROPOSTA e che non sono state rifiutate tramite RIFIUTO.} \\
\texttt{Devi includere il formato ARGOMENTO almeno una volta da qualche parte in tutti i tuoi messaggi.} \\
\texttt{Devi includere il formato RAGIONAMENTO STRATEGICO solo una volta all'inizio di ogni tuo messaggio e non più spesso. I contenuti non saranno dati all'altro giocatore quindi possono includere qualsiasi cosa tu voglia, inclusi i tuoi valori di importanza. Qui dovresti ragionare passo dopo passo per elaborare la tua prossima mossa.} \\
\\ 
\texttt{Ora riceverai il primo messaggio dell'altro giocatore.} \\
\\ 
\texttt{ARGOMENTAZIONE: \{"La mia proposta iniziale sarà basata sugli oggetti con il più alto rapporto tra importanza e sforzo. Ciò garantisce che il set iniziale sia ottimizzato per massimizzare l'importanza totale senza superare il limite di sforzo. Sono aperto a negoziare e adattare la mia proposta in base alle indicazioni dell'altro partecipante."\}} \\
\\ 
\texttt{PROPOSTA: \{'C54', 'A04', 'B41', 'A10'\}} \\
            }
        }
    }
    & \\ \\

    \theutterance \stepcounter{utterance}  
    & & & \multicolumn{4}{p{0.6\linewidth}}{
        \cellcolor[rgb]{0.9,0.9,0.9}{
            \makecell[{{p{\linewidth}}}]{
                \texttt{\tiny{[GM$\langle$P2]}}
                \texttt{RAGIONAMENTO STRATEGICO: \{'Sto analizzando la proposta iniziale dell\textbackslash{}'altro partecipante, che include gli oggetti C54, A04, B41 e A10. Il loro sforzo totale è 313 + 203 + 720 + 437 = 1673, che è al di sotto del limite di 3386. La mia strategia sarà quella di valutare l\textbackslash{}'importanza di questa proposta in base ai miei valori di importanza personali e di cercare di massimizzare l\textbackslash{}'importanza totale senza superare il limite di sforzo. Considererò anche la possibilità di sostituire alcuni oggetti con altri che potrebbero avere un rapporto più alto tra importanza e sforzo.'\}} \\
\\ 
\texttt{ARGOMENTAZIONE: \{"La proposta iniziale sembra ragionevole, ma potrebbe essere migliorata. Considererò l\textbackslash{}'aggiunta di altri oggetti che potrebbero aumentare l\textbackslash{}'importanza totale senza superare il limite di sforzo. Sono disposto a negoziare e a trovare un compromesso che soddisfi entrambi i partecipanti."\}} \\
\\ 
\texttt{PROPOSTA: \{'C54', 'A04', 'B41', 'A10', 'C01'\}} \\
            }
        }
    }
     \\ \\

    \theutterance \stepcounter{utterance}  
    & & \multicolumn{4}{p{0.6\linewidth}}{
        \cellcolor[rgb]{0.9,0.9,0.9}{
            \makecell[{{p{\linewidth}}}]{
                \texttt{\tiny{[P1$\langle$GM]}}
                \texttt{ARGOMENTAZIONE: \{"La proposta iniziale sembra ragionevole, ma potrebbe essere migliorata. Considererò l\textbackslash{}'aggiunta di altri oggetti che potrebbero aumentare l\textbackslash{}'importanza totale senza superare il limite di sforzo. Sono disposto a negoziare e a trovare un compromesso che soddisfi entrambi i partecipanti."\}} \\
\\ 
\texttt{PROPOSTA: \{'C54', 'A04', 'B41', 'A10', 'C01'\}} \\
            }
        }
    }
    & \\ \\

    \theutterance \stepcounter{utterance}  
    & \multicolumn{4}{p{0.6\linewidth}}{
        \cellcolor[rgb]{0.9,0.9,0.9}{
            \makecell[{{p{\linewidth}}}]{
                \texttt{\tiny{[P1$\rangle$GM]}}
                \texttt{RAGIONAMENTO STRATEGICO: \{"L'aggiunta di 'C01' alla proposta precedente potrebbe essere vantaggiosa, ma è necessario verificare se il set proposto rimane entro il limite di sforzo. Calcolerò lo sforzo totale del set proposto e valuterò se ci sono altri oggetti che potrebbero essere aggiunti o rimossi per migliorare l'importanza totale."\}} \\
\\ 
\texttt{ARGOMENTAZIONE: \{"La proposta aggiornata include 'C01', che ha un elevato valore di importanza. Tuttavia, è fondamentale verificare che il set proposto non superi il limite di sforzo. Sono disposto a valutare altre opzioni e a trovare un equilibrio tra importanza e sforzo."\}} \\
\\ 
\texttt{PROPOSTA: \{'C54', 'A04', 'B41', 'A10', 'C01', 'A87'\}} \\
            }
        }
    }
    & & \\ \\

    \theutterance \stepcounter{utterance}  
    & & \multicolumn{4}{p{0.6\linewidth}}{
        \cellcolor[rgb]{0.9,0.9,0.9}{
            \makecell[{{p{\linewidth}}}]{
                \texttt{\tiny{[GM$\rangle$P2]}}
                \texttt{ARGOMENTAZIONE: \{"La proposta aggiornata include 'C01', che ha un elevato valore di importanza. Tuttavia, è fondamentale verificare che il set proposto non superi il limite di sforzo. Sono disposto a valutare altre opzioni e a trovare un equilibrio tra importanza e sforzo."\}} \\
\\ 
\texttt{PROPOSTA: \{'C54', 'A04', 'B41', 'A10', 'C01', 'A87'\}} \\
            }
        }
    }
    & \\ \\

    \theutterance \stepcounter{utterance}  
    & & & \multicolumn{4}{p{0.6\linewidth}}{
        \cellcolor[rgb]{0.9,0.9,0.9}{
            \makecell[{{p{\linewidth}}}]{
                \texttt{\tiny{[GM$\langle$P2]}}
                \texttt{RAGIONAMENTO STRATEGICO: \{'La nuova proposta include A87, che ha un sforzo molto basso (4) e potrebbe essere un ottimo candidato per essere inclusa nel set finale. Tuttavia, devo considerare l\textbackslash{}'impatto della sua aggiunta sul set completo e valutare se il set proposto supera il limite di sforzo. Il sforzo totale del set proposto è 313 + 203 + 720 + 437 + 539 + 4 = 2216, che è ancora al di sotto del limite di 3386. La mia strategia sarà quella di continuare a valutare le opzioni e cercare di massimizzare l\textbackslash{}'importanza totale del set.'\}} \\
\\ 
\texttt{ARGOMENTAZIONE: \{"La proposta aggiornata sembra promettente, ma potrebbe essere ulteriormente migliorata. Considererò l\textbackslash{}'aggiunta di altri oggetti che potrebbero aumentare l\textbackslash{}'importanza totale senza superare il limite di sforzo. Sono disposto a negoziare e a trovare un compromesso che soddisfi entrambi i partecipanti."\}} \\
\\ 
\texttt{RIFIUTO: \{'C54', 'A04', 'B41', 'A10', 'C01', 'A87'\} } \\
\\ 
\texttt{PROPOSTA: \{'C54', 'A04', 'B41', 'A10', 'C01', 'A87', 'B33'\}} \\
            }
        }
    }
     \\ \\

    \theutterance \stepcounter{utterance}  
    & & \multicolumn{4}{p{0.6\linewidth}}{
        \cellcolor[rgb]{0.9,0.9,0.9}{
            \makecell[{{p{\linewidth}}}]{
                \texttt{\tiny{[P1$\langle$GM]}}
                \texttt{ARGOMENTAZIONE: \{"La proposta aggiornata sembra promettente, ma potrebbe essere ulteriormente migliorata. Considererò l\textbackslash{}'aggiunta di altri oggetti che potrebbero aumentare l\textbackslash{}'importanza totale senza superare il limite di sforzo. Sono disposto a negoziare e a trovare un compromesso che soddisfi entrambi i partecipanti."\}} \\
\\ 
\texttt{RIFIUTO: \{'C54', 'A04', 'B41', 'A10', 'C01', 'A87'\} } \\
\\ 
\texttt{PROPOSTA: \{'C54', 'A04', 'B41', 'A10', 'C01', 'A87', 'B33'\}} \\
            }
        }
    }
    & \\ \\

    \theutterance \stepcounter{utterance}  
    & \multicolumn{4}{p{0.6\linewidth}}{
        \cellcolor[rgb]{0.9,0.9,0.9}{
            \makecell[{{p{\linewidth}}}]{
                \texttt{\tiny{[P1$\rangle$GM]}}
                \texttt{RAGIONAMENTO STRATEGICO: \{"Il rifiuto della proposta precedente e l'aggiunta di 'B33' alla nuova proposta indicano che l'altro partecipante sta cercando di massimizzare l'importanza totale. Tuttavia, devo verificare se il set proposto rimane entro il limite di sforzo. Calcolerò lo sforzo totale del set proposto e valuterò se l'aggiunta di 'B33' è vantaggiosa o se ci sono altri oggetti che potrebbero essere aggiunti o rimossi per migliorare l'importanza totale."\}} \\
\\ 
\texttt{ARGOMENTAZIONE: \{"La proposta aggiornata include 'B33', che ha un basso valore di importanza. Sono preoccupato che l'aggiunta di questo oggetto possa non essere la scelta più ottimale. Sono disposto a valutare altre opzioni e a trovare un equilibrio tra importanza e sforzo. Potrebbe essere necessario rimuovere 'B33' e cercare alternative più vantaggiose."\}} \\
\\ 
\texttt{RIFIUTO: \{'C54', 'A04', 'B41', 'A10', 'C01', 'A87', 'B33'\}} \\
            }
        }
    }
    & & \\ \\

    \theutterance \stepcounter{utterance}  
    & & \multicolumn{4}{p{0.6\linewidth}}{
        \cellcolor[rgb]{0.9,0.9,0.9}{
            \makecell[{{p{\linewidth}}}]{
                \texttt{\tiny{[GM$\rangle$P2]}}
                \texttt{ARGOMENTAZIONE: \{"La proposta aggiornata include 'B33', che ha un basso valore di importanza. Sono preoccupato che l'aggiunta di questo oggetto possa non essere la scelta più ottimale. Sono disposto a valutare altre opzioni e a trovare un equilibrio tra importanza e sforzo. Potrebbe essere necessario rimuovere 'B33' e cercare alternative più vantaggiose."\}} \\
\\ 
\texttt{RIFIUTO: \{'C54', 'A04', 'B41', 'A10', 'C01', 'A87', 'B33'\}} \\
            }
        }
    }
    & \\ \\

    \theutterance \stepcounter{utterance}  
    & & & \multicolumn{4}{p{0.6\linewidth}}{
        \cellcolor[rgb]{0.9,0.9,0.9}{
            \makecell[{{p{\linewidth}}}]{
                \texttt{\tiny{[GM$\langle$P2]}}
                \texttt{RAGIONAMENTO STRATEGICO: \{'La proposta precedente è stata rifiutata e l\textbackslash{}'altro partecipante ha espresso preoccupazione sull\textbackslash{}'aggiunta di B33. Considererò di sostituire B33 con un altro oggetto che abbia un rapporto più alto tra importanza e sforzo. Un\textbackslash{}'opzione potrebbe essere A77, che ha un elevato valore di importanza e un sforzo moderate (685). Tuttavia, devo verificare se l\textbackslash{}'aggiunta di A77 supera il limite di sforzo.'\}} \\
\\ 
\texttt{ARGOMENTAZIONE: \{"Sono d'accordo che B33 potrebbe non essere la scelta più ottimale. Considererò di sostituire questo oggetto con un altro che abbia un miglior rapporto tra importanza e sforzo. Sono disposto a valutare altre opzioni e a trovare un compromesso che soddisfi entrambi i partecipanti."\}} \\
\\ 
\texttt{PROPOSTA: \{'C54', 'A04', 'B41', 'A10', 'C01', 'A87', 'A77'\}} \\
            }
        }
    }
     \\ \\

    \theutterance \stepcounter{utterance}  
    & & \multicolumn{4}{p{0.6\linewidth}}{
        \cellcolor[rgb]{0.9,0.9,0.9}{
            \makecell[{{p{\linewidth}}}]{
                \texttt{\tiny{[P1$\langle$GM]}}
                \texttt{ARGOMENTAZIONE: \{"Sono d'accordo che B33 potrebbe non essere la scelta più ottimale. Considererò di sostituire questo oggetto con un altro che abbia un miglior rapporto tra importanza e sforzo. Sono disposto a valutare altre opzioni e a trovare un compromesso che soddisfi entrambi i partecipanti."\}} \\
\\ 
\texttt{PROPOSTA: \{'C54', 'A04', 'B41', 'A10', 'C01', 'A87', 'A77'\}} \\
            }
        }
    }
    & \\ \\

    \theutterance \stepcounter{utterance}  
    & \multicolumn{4}{p{0.6\linewidth}}{
        \cellcolor[rgb]{0.9,0.9,0.9}{
            \makecell[{{p{\linewidth}}}]{
                \texttt{\tiny{[P1$\rangle$GM]}}
                \texttt{RAGIONAMENTO STRATEGICO: \{"La sostituzione di 'B33' con 'A77' nella proposta rappresenta un significativo miglioramento, dato che 'A77' ha un alto valore di importanza e un rapporto tra importanza e sforzo più equilibrato. Devo calcolare lo sforzo totale del set proposto per assicurarmi che non superi il limite di sforzo. Inoltre, valuterò se questa combinazione di oggetti massimizza l'importanza totale senza eccedere il limite di sforzo.")} \\
\\ 
\texttt{ARGOMENTAZIONE: \{"La proposta aggiornata con 'A77' al posto di 'B33' sembra molto più equilibrata e vantaggiosa. 'A77' contribuisce notevolmente all'importanza totale del set. Sono favorevole a questa proposta, poiché sembra raggiungere un buon equilibrio tra importanza e sforzo, e sono disposto ad accettarla come nostro set condiviso."\}} \\
\\ 
\texttt{ACCORDO: \{'C54', 'A04', 'B41', 'A10', 'C01', 'A87', 'A77'\}} \\
            }
        }
    }
    & & \\ \\

    \theutterance \stepcounter{utterance}  
    & & \multicolumn{4}{p{0.6\linewidth}}{
        \cellcolor[rgb]{0.9,0.9,0.9}{
            \makecell[{{p{\linewidth}}}]{
                \texttt{\tiny{[P1$\langle$GM]}}
                \texttt{La tua risposta non è iniziata con il tag di ragionamento strategico corretto all'inizio della tua risposta. Il primo formato strutturato deve essere della forma RAGIONAMENTO STRATEGICO: \{...\}. Riprova.} \\
            }
        }
    }
    & \\ \\

    \theutterance \stepcounter{utterance}  
    & \multicolumn{4}{p{0.6\linewidth}}{
        \cellcolor[rgb]{0.9,0.9,0.9}{
            \makecell[{{p{\linewidth}}}]{
                \texttt{\tiny{[P1$\rangle$GM]}}
                \texttt{RAGIONAMENTO STRATEGICO: \{"La sostituzione di 'B33' con 'A77' nella proposta rappresenta un significativo miglioramento, dato che 'A77' ha un alto valore di importanza e un rapporto tra importanza e sforzo più equilibrato. Devo calcolare lo sforzo totale del set proposto per assicurarmi che non superi il limite di sforzo. Inoltre, valuterò se questa combinazione di oggetti massimizza l'importanza totale senza eccedere il limite di sforzo. Considererò anche la strategia dell'altro partecipante e la sua possibile volontà di negoziare ulteriormente."\}} \\
\\ 
\texttt{ARGOMENTAZIONE: \{"La proposta aggiornata con 'A77' al posto di 'B33' sembra molto più equilibrata e vantaggiosa. 'A77' contribuisce notevolmente all'importanza totale del set. Sono favorevole a questa proposta, poiché sembra raggiungere un buon equilibrio tra importanza e sforzo, e sono disposto ad accettarla come nostro set condiviso."\}} \\
\\ 
\texttt{ACCORDO: \{'C54', 'A04', 'B41', 'A10', 'C01', 'A87', 'A77'\}} \\
            }
        }
    }
    & & \\ \\

    \theutterance \stepcounter{utterance}  
    & & & \multicolumn{2}{p{0.3\linewidth}}{
        \cellcolor[rgb]{0.9,0.9,0.9}{
            \makecell[{{p{\linewidth}}}]{
                \texttt{\tiny{[GM$|$GM]}}
                \texttt{game successful} \\
            }
        }
    }
    & & \\ \\

    \theutterance \stepcounter{utterance}  
    & & & \multicolumn{2}{p{0.3\linewidth}}{
        \cellcolor[rgb]{0.9,0.9,0.9}{
            \makecell[{{p{\linewidth}}}]{
                \texttt{\tiny{[GM$|$GM]}}
                \texttt{end game} \\
            }
        }
    }
    & & \\ \\

\end{supertabular}
}

\end{document}
